% assets/a3-alloc-chain.svg — 제어기에서 운동모델까지, 배분 사슬 전체
%
%   bash _tools/tikz2svg.sh assets/src/a3-alloc-chain.tex assets/a3-alloc-chain.svg
%
% 이 그림이 말하는 것
%   제어기가 낸 요구 렌치 tau_c = [X; Y; N] 이
%   (1) pseudo-inverse -> (2) null space 보정 -> (3) 극좌표 -> (4) 한계 처리
%   를 거쳐 모터 추력 둘과 노즐 각도 둘이 되고,
%   그것이 같은 행렬 B 로 **되돌아가** 실제 렌치 tau_a 가 된 다음 운동모델로 들어간다.
%   요구와 실제의 차 e = tau_a - tau_c 가 실측표의 '못 낸 X, Y, N' 이다.
%
% 배치 — 세 줄 사슬 (좌->우, 꺾어 내려온다)
%   A 줄 y = 0    제어기 -> (1) -> (2) -> (3)
%   B 줄 y = -52  (4) -> (5) -> 운동모델 -> 항적
%   C 줄 y = -92  요구와 실제의 비교
%   폭 186 mm 에 맞춘다 (PDF 본문 폭). 상자 36 mm, 간격 14 mm

\documentclass[border=6pt]{standalone}
% capstone-style.tex — 이 볼트의 TikZ 그림이 공유하는 스타일
%
% 쓰는 법:  % capstone-style.tex — 이 볼트의 TikZ 그림이 공유하는 스타일
%
% 쓰는 법:  % capstone-style.tex — 이 볼트의 TikZ 그림이 공유하는 스타일
%
% 쓰는 법:  \input{capstone-style}   (assets/src/ 안에서)
% 빌드:     bash _tools/tikz2svg.sh assets/src/w02-node-topic.tex assets/w02-node-topic.svg
%
% 왜 TikZ 인가
%   손으로 SVG 좌표를 쓰면 화살촉·선 굵기·글자 크기가 그림마다 미묘하게 달라진다.
%   그것이 "자료가 어설퍼 보인다"의 정체다. 스타일을 한 번만 정하고 상대좌표로
%   배치하면 좌표를 고쳐도 그림이 무너지지 않는다.
%   수식은 본문과 같은 TeX 글꼴로 나온다.
%
% 한글
%   ko.TeX(DVI 경로)를 쓴다. XeLaTeX 은 TikZ 그래픽을 PDF special 로 내보내는데
%   dvisvgm 이 그것을 받으려면 Ghostscript < 10.01 이 필요하다. 이 PC 에는
%   10.07 이 깔려 있어 그 경로가 막힌다. latex -> DVI -> dvisvgm 이 유일하게
%   동작하는 조합이다.

\usepackage[hangul]{kotex}
\usepackage{tikz}
\usepackage{amsmath}
\usepackage{xcolor}
\usetikzlibrary{arrows.meta, positioning, calc, fit, backgrounds, decorations.pathmorphing}

% ---- 색 --------------------------------------------------------------------
% 강조는 하나만. 나머지는 검정. 흑백 인쇄에서도 읽혀야 한다.
\definecolor{ink}{HTML}{1A1A1A}
\definecolor{soft}{HTML}{666666}
\definecolor{accent}{HTML}{7C3AED}
\definecolor{warn}{HTML}{B91C1C}

% 계층 색 — Simulink 모델의 블록 색과 같은 규칙을 쓴다
\definecolor{guid}{HTML}{1D4ED8}    % 유도
\definecolor{ctrl}{HTML}{EA580C}    % 제어
\definecolor{thr}{HTML}{D97706}     % 추진기
\definecolor{plant}{HTML}{16A34A}   % 운동모델
\definecolor{nav}{HTML}{7C3AED}     % 항법 · 신호처리
\definecolor{pathc}{HTML}{0D9488}   % 계획 경로 · 항적
\definecolor{errc}{HTML}{C2410C}    % 오차
\definecolor{flow}{HTML}{0369A1}    % 조류 · 바람

% ---- 선과 화살촉 -----------------------------------------------------------
\tikzset{
  sig/.style   = {ink, line width=0.5pt, -{Stealth[length=4pt, width=3pt]}},
  wire/.style  = {ink, line width=0.5pt},
  fb/.style    = {ink, line width=0.5pt, densely dashed,
                  -{Stealth[length=4pt, width=3pt]}},
  asig/.style  = {accent, line width=0.5pt, densely dashed,
                  -{Stealth[length=4pt, width=3pt]}},
  % 블록 — 흰 바탕, 직각, 검은 테두리. 색은 테두리에만 준다
  blk/.style   = {draw=ink, line width=0.55pt, fill=white,
                  minimum width=18mm, minimum height=8.5mm, inner sep=2pt, font=\small},
  gblk/.style  = {blk, draw=guid},
  cblk/.style  = {blk, draw=ctrl},
  tblk/.style  = {blk, draw=thr},
  pblk/.style  = {blk, draw=plant},
  nblk/.style  = {blk, draw=nav},
  % 합산점 · 분기점
  sum/.style   = {draw=ink, line width=0.55pt, fill=white, circle,
                  inner sep=0pt, minimum size=4.6mm, font=\scriptsize},
  dot/.style   = {ink, fill=ink, circle, inner sep=0pt, minimum size=1.5mm},
  % 글자
  sname/.style = {font=\small, inner sep=1.5pt},
  note/.style  = {font=\scriptsize, soft, inner sep=1.5pt},
  wnote/.style = {font=\scriptsize, warn, inner sep=1.5pt},
  anote/.style = {font=\scriptsize, accent, inner sep=1.5pt},
  axis/.style  = {ink, line width=0.5pt},
  tick/.style  = {font=\scriptsize, soft},
  % 묶음 상자 — 서브시스템 표시
  grp/.style   = {draw=soft, line width=0.4pt, densely dotted, rounded corners=1mm,
                  inner sep=3mm},
}
   (assets/src/ 안에서)
% 빌드:     bash _tools/tikz2svg.sh assets/src/w02-node-topic.tex assets/w02-node-topic.svg
%
% 왜 TikZ 인가
%   손으로 SVG 좌표를 쓰면 화살촉·선 굵기·글자 크기가 그림마다 미묘하게 달라진다.
%   그것이 "자료가 어설퍼 보인다"의 정체다. 스타일을 한 번만 정하고 상대좌표로
%   배치하면 좌표를 고쳐도 그림이 무너지지 않는다.
%   수식은 본문과 같은 TeX 글꼴로 나온다.
%
% 한글
%   ko.TeX(DVI 경로)를 쓴다. XeLaTeX 은 TikZ 그래픽을 PDF special 로 내보내는데
%   dvisvgm 이 그것을 받으려면 Ghostscript < 10.01 이 필요하다. 이 PC 에는
%   10.07 이 깔려 있어 그 경로가 막힌다. latex -> DVI -> dvisvgm 이 유일하게
%   동작하는 조합이다.

\usepackage[hangul]{kotex}
\usepackage{tikz}
\usepackage{amsmath}
\usepackage{xcolor}
\usetikzlibrary{arrows.meta, positioning, calc, fit, backgrounds, decorations.pathmorphing}

% ---- 색 --------------------------------------------------------------------
% 강조는 하나만. 나머지는 검정. 흑백 인쇄에서도 읽혀야 한다.
\definecolor{ink}{HTML}{1A1A1A}
\definecolor{soft}{HTML}{666666}
\definecolor{accent}{HTML}{7C3AED}
\definecolor{warn}{HTML}{B91C1C}

% 계층 색 — Simulink 모델의 블록 색과 같은 규칙을 쓴다
\definecolor{guid}{HTML}{1D4ED8}    % 유도
\definecolor{ctrl}{HTML}{EA580C}    % 제어
\definecolor{thr}{HTML}{D97706}     % 추진기
\definecolor{plant}{HTML}{16A34A}   % 운동모델
\definecolor{nav}{HTML}{7C3AED}     % 항법 · 신호처리
\definecolor{pathc}{HTML}{0D9488}   % 계획 경로 · 항적
\definecolor{errc}{HTML}{C2410C}    % 오차
\definecolor{flow}{HTML}{0369A1}    % 조류 · 바람

% ---- 선과 화살촉 -----------------------------------------------------------
\tikzset{
  sig/.style   = {ink, line width=0.5pt, -{Stealth[length=4pt, width=3pt]}},
  wire/.style  = {ink, line width=0.5pt},
  fb/.style    = {ink, line width=0.5pt, densely dashed,
                  -{Stealth[length=4pt, width=3pt]}},
  asig/.style  = {accent, line width=0.5pt, densely dashed,
                  -{Stealth[length=4pt, width=3pt]}},
  % 블록 — 흰 바탕, 직각, 검은 테두리. 색은 테두리에만 준다
  blk/.style   = {draw=ink, line width=0.55pt, fill=white,
                  minimum width=18mm, minimum height=8.5mm, inner sep=2pt, font=\small},
  gblk/.style  = {blk, draw=guid},
  cblk/.style  = {blk, draw=ctrl},
  tblk/.style  = {blk, draw=thr},
  pblk/.style  = {blk, draw=plant},
  nblk/.style  = {blk, draw=nav},
  % 합산점 · 분기점
  sum/.style   = {draw=ink, line width=0.55pt, fill=white, circle,
                  inner sep=0pt, minimum size=4.6mm, font=\scriptsize},
  dot/.style   = {ink, fill=ink, circle, inner sep=0pt, minimum size=1.5mm},
  % 글자
  sname/.style = {font=\small, inner sep=1.5pt},
  note/.style  = {font=\scriptsize, soft, inner sep=1.5pt},
  wnote/.style = {font=\scriptsize, warn, inner sep=1.5pt},
  anote/.style = {font=\scriptsize, accent, inner sep=1.5pt},
  axis/.style  = {ink, line width=0.5pt},
  tick/.style  = {font=\scriptsize, soft},
  % 묶음 상자 — 서브시스템 표시
  grp/.style   = {draw=soft, line width=0.4pt, densely dotted, rounded corners=1mm,
                  inner sep=3mm},
}
   (assets/src/ 안에서)
% 빌드:     bash _tools/tikz2svg.sh assets/src/w02-node-topic.tex assets/w02-node-topic.svg
%
% 왜 TikZ 인가
%   손으로 SVG 좌표를 쓰면 화살촉·선 굵기·글자 크기가 그림마다 미묘하게 달라진다.
%   그것이 "자료가 어설퍼 보인다"의 정체다. 스타일을 한 번만 정하고 상대좌표로
%   배치하면 좌표를 고쳐도 그림이 무너지지 않는다.
%   수식은 본문과 같은 TeX 글꼴로 나온다.
%
% 한글
%   ko.TeX(DVI 경로)를 쓴다. XeLaTeX 은 TikZ 그래픽을 PDF special 로 내보내는데
%   dvisvgm 이 그것을 받으려면 Ghostscript < 10.01 이 필요하다. 이 PC 에는
%   10.07 이 깔려 있어 그 경로가 막힌다. latex -> DVI -> dvisvgm 이 유일하게
%   동작하는 조합이다.

\usepackage[hangul]{kotex}
\usepackage{tikz}
\usepackage{amsmath}
\usepackage{xcolor}
\usetikzlibrary{arrows.meta, positioning, calc, fit, backgrounds, decorations.pathmorphing}

% ---- 색 --------------------------------------------------------------------
% 강조는 하나만. 나머지는 검정. 흑백 인쇄에서도 읽혀야 한다.
\definecolor{ink}{HTML}{1A1A1A}
\definecolor{soft}{HTML}{666666}
\definecolor{accent}{HTML}{7C3AED}
\definecolor{warn}{HTML}{B91C1C}

% 계층 색 — Simulink 모델의 블록 색과 같은 규칙을 쓴다
\definecolor{guid}{HTML}{1D4ED8}    % 유도
\definecolor{ctrl}{HTML}{EA580C}    % 제어
\definecolor{thr}{HTML}{D97706}     % 추진기
\definecolor{plant}{HTML}{16A34A}   % 운동모델
\definecolor{nav}{HTML}{7C3AED}     % 항법 · 신호처리
\definecolor{pathc}{HTML}{0D9488}   % 계획 경로 · 항적
\definecolor{errc}{HTML}{C2410C}    % 오차
\definecolor{flow}{HTML}{0369A1}    % 조류 · 바람

% ---- 선과 화살촉 -----------------------------------------------------------
\tikzset{
  sig/.style   = {ink, line width=0.5pt, -{Stealth[length=4pt, width=3pt]}},
  wire/.style  = {ink, line width=0.5pt},
  fb/.style    = {ink, line width=0.5pt, densely dashed,
                  -{Stealth[length=4pt, width=3pt]}},
  asig/.style  = {accent, line width=0.5pt, densely dashed,
                  -{Stealth[length=4pt, width=3pt]}},
  % 블록 — 흰 바탕, 직각, 검은 테두리. 색은 테두리에만 준다
  blk/.style   = {draw=ink, line width=0.55pt, fill=white,
                  minimum width=18mm, minimum height=8.5mm, inner sep=2pt, font=\small},
  gblk/.style  = {blk, draw=guid},
  cblk/.style  = {blk, draw=ctrl},
  tblk/.style  = {blk, draw=thr},
  pblk/.style  = {blk, draw=plant},
  nblk/.style  = {blk, draw=nav},
  % 합산점 · 분기점
  sum/.style   = {draw=ink, line width=0.55pt, fill=white, circle,
                  inner sep=0pt, minimum size=4.6mm, font=\scriptsize},
  dot/.style   = {ink, fill=ink, circle, inner sep=0pt, minimum size=1.5mm},
  % 글자
  sname/.style = {font=\small, inner sep=1.5pt},
  note/.style  = {font=\scriptsize, soft, inner sep=1.5pt},
  wnote/.style = {font=\scriptsize, warn, inner sep=1.5pt},
  anote/.style = {font=\scriptsize, accent, inner sep=1.5pt},
  axis/.style  = {ink, line width=0.5pt},
  tick/.style  = {font=\scriptsize, soft},
  % 묶음 상자 — 서브시스템 표시
  grp/.style   = {draw=soft, line width=0.4pt, densely dotted, rounded corners=1mm,
                  inner sep=3mm},
}


\begin{document}
\begin{tikzpicture}[x=1mm, y=1mm]

\tikzset{
  bx/.style  = {rectangle, rounded corners=1pt, draw, line width=0.7pt,
                align=center, font=\scriptsize, inner sep=2.5pt,
                minimum width=36mm, minimum height=18mm},
  sbx/.style = {rectangle, rounded corners=1pt, draw, line width=0.7pt,
                align=center, font=\scriptsize, inner sep=2.5pt,
                minimum width=36mm, minimum height=12mm},
  ar/.style  = {-{Stealth[length=4pt,width=3pt]}, ink, line width=0.7pt},
  dar/.style = {-{Stealth[length=4pt,width=3pt]}, soft, line width=0.6pt,
                densely dashed},
  lb/.style  = {font=\scriptsize, ink, anchor=south, inner sep=1pt},
}

% ===========================================================================
% A 줄 — 제어기와 배분 세 걸음
% ===========================================================================
\node[bx, draw=ctrl, fill=ctrl!6] (a1) at (18,0)
     {\textbf{제어기} (W06 과 동일) \\[2pt] 위치$\cdot$선수각 오차 \\ $\longrightarrow$ 요구 렌치};

\node[bx, draw=guid, fill=guid!6] (a2) at (68,0)
     {\textbf{\textcircled{\scriptsize 1} pseudo-inverse} \\[2pt]
      $\mathbf{u}_e = \mathbf{B}^{+}\boldsymbol\tau_c$ \\[2pt]
      크기가 가장 작은 해};

\node[bx, draw=guid, fill=guid!6] (a3) at (118,0)
     {\textbf{\textcircled{\scriptsize 2} null space 보정} \\[2pt]
      $\mathbf{u}_e \leftarrow \mathbf{u}_e + \lambda\,\mathbf{n}$ \\[2pt]
      각도를 쐐기 밖으로};

\node[bx, draw=thr, fill=thr!6] (a4) at (168,0)
     {\textbf{\textcircled{\scriptsize 3} 극좌표 변환} \\[2pt]
      $T=\sqrt{F_x^2+F_y^2}$ \\[1pt]
      $\delta=\operatorname{atan2}(F_y, F_x)$};

\draw[ar] (a1) -- (a2);
\draw[ar] (a2) -- (a3);
\draw[ar] (a3) -- (a4);

\node[lb] at (43,10.5)  {$\boldsymbol\tau_c = [\,X,\ Y,\ N\,]^{\top}$ \ (3$\times$1)};
\node[lb] at (93,10.5)  {$\mathbf{u}_e = [\,F_{x,L}, F_{y,L}, F_{x,R}, F_{y,R}\,]^{\top}$};
\node[lb] at (143,10.5) {보정된 $\mathbf{u}_e$};

% ===========================================================================
% 꺾어 내려오는 선
% ===========================================================================
\draw[ar] (a4.south) -- (168,-28) -- (18,-28) -- (18,-43);
\node[lb, soft] at (93,-27) {요구 추력 $T_L, T_R$ \quad 요구 각도 $\delta_L, \delta_R$};

% ===========================================================================
% B 줄 — 한계 처리, 실제 렌치, 운동모델
% ===========================================================================
\node[bx, draw=warn, fill=warn!6] (b1) at (18,-52)
     {\textbf{\textcircled{\scriptsize 4} 한계 처리} \\[2pt]
      각도 $\pm45^\circ$ $\cdot$ 변화율 $573^\circ$/s \\[1pt]
      추력 $\pm500$ N};

\node[bx, draw=plant, fill=plant!6] (b2) at (68,-52)
     {\textbf{\textcircled{\scriptsize 5} 실제 렌치} \\[2pt]
      $\boldsymbol\tau_a = \mathbf{B}\,\mathbf{R}(\boldsymbol\delta_a)\,\mathbf{T}_a$ \\[2pt]
      같은 $\mathbf{B}$ 를 되짚는다};

\node[bx, draw=plant, fill=plant!6] (b3) at (118,-52)
     {\textbf{운동모델} (3 자유도) \\[2pt]
      $\mathbf{M}\dot{\boldsymbol\nu} + \mathbf{C}(\boldsymbol\nu)\boldsymbol\nu
       + \mathbf{D}\boldsymbol\nu = \boldsymbol\tau_a$};

\node[bx, draw=nav, fill=nav!6] (b4) at (168,-52)
     {\textbf{항적 $\cdot$ 상태} \\[2pt]
      $x,\ y,\ \psi,\ u,\ v,\ r$};

\draw[ar] (b1) -- (b2);
\draw[ar] (b2) -- (b3);
\draw[ar] (b3) -- (b4);

\node[lb] at (43,-41.5)  {실제 $\mathbf{T}_a$, $\boldsymbol\delta_a$};
\node[lb] at (93,-41.5)  {$\boldsymbol\tau_a = [\,X_a,\ Y_a,\ N_a\,]^{\top}$};
\node[lb] at (143,-41.5) {$\boldsymbol\nu = [\,u,\ v,\ r\,]^{\top}$};

% ===========================================================================
% C 줄 — 요구와 실제를 맞대어 본다
% ===========================================================================
\node[sbx, draw=soft, dashed, fill=white] (c1) at (18,-92)
     {요구 $\boldsymbol\tau_c$ \\ (\textcircled{\scriptsize 1} 의 입력 그대로)};

\node[circle, draw=ink, line width=0.7pt, minimum size=5mm, inner sep=0pt,
      font=\scriptsize] (sum) at (68,-92) {$-$};

\node[sbx, draw=errc, fill=errc!6, minimum width=70mm] (c2) at (143,-92)
     {\textbf{못 낸 힘} \quad
      $\mathbf{e} = \boldsymbol\tau_a - \boldsymbol\tau_c$ \\[2pt]
      실측표의 `못 낸 $X\cdot Y\cdot N$' 가 이 값이다};

\draw[dar] (c1) -- (sum);
\draw[dar] (sum) -- (c2);
\draw[dar] (b2.south) -- (68,-74) -- (sum);
\node[font=\scriptsize, soft, anchor=west, inner sep=1pt] at (70,-70)
     {실제 $\boldsymbol\tau_a$};

% ===========================================================================
% 아래 띠 — 읽는 법
% ===========================================================================
\draw[wire, soft] (0,-108) -- (186,-108);

\node[anchor=north west, font=\scriptsize, text width=88mm] at (0,-113)
     {\textbf{사슬에서 꼭 볼 곳 둘} \\[4pt]
      \textbf{(1)} \textcircled{\scriptsize 1}$\sim$\textcircled{\scriptsize 3} 은
      $\mathbf{B}$ 를 \textbf{거꾸로} 푸는 과정이고,
      \textcircled{\scriptsize 5} 는 같은 $\mathbf{B}$ 를 \textbf{바로} 쓰는 과정이다.
      행렬 하나가 배분과 검산에 함께 쓰인다 \\[4pt]
      \textbf{(2)} 한계에 물리지 않으면
      $\boldsymbol\tau_a = \boldsymbol\tau_c$ 가 되어 $\mathbf{e}=\mathbf{0}$ 이다.
      기준 환경 실측에서 한계 밖 구간의 차이는 $3\times10^{-14} \sim 1.4\times10^{-12}$
      — 수치 오차 수준};

\node[anchor=north west, font=\scriptsize, text width=88mm] at (98,-113)
     {\textbf{코드에서 어디에 있는가} \\[4pt]
      \textcircled{\scriptsize 1}$\sim$\textcircled{\scriptsize 4}
      \quad \texttt{parts/azAllocFcn.m} \ (AzAlloc 블록) \\[2pt]
      \textcircled{\scriptsize 5}
      \quad \texttt{parts/wrenchFcn.m} \ (Wrench 블록) \\[2pt]
      운동모델 \quad \texttt{A3\_*.slx} 의 Plant 서브시스템 \\[4pt]
      두 함수가 같은 $\mathbf{B}$ 를 각자 한 줄로 적어 둔다 —
      기하를 고치면 \textbf{두 곳을 함께} 고쳐야 한다};

\node[anchor=north west, font=\scriptsize, warn, text width=186mm] at (0,-146)
     {\textcircled{\scriptsize 2} 는 Sørdalen(1997)과 Sarda 등(2016)이 쓴 null space 이용법을
      이 배치(추진기 둘)에 맞춰 닫힌식으로 줄인 것이다 — 반복 최적화를 돌리지 않는다};

\end{tikzpicture}
\end{document}
